% Appendix "comp" (writer: oodcomp). Proofs from research/tracks/context/notes-final.md §3 (Thm 3.1, Prop 3.2c, Prop 3.9,
% Thm 3.5, Lemma 3.11a, Thm 3.11); computations c1, c5, c6, c8, c9, c11 in research/tracks/context/checks/.
\section{Proofs for Section~\ref{sec:comp}}\label{app:comp}

\begin{proof}[Proof of \cref{thm:comp:timebounded}(a)]
$\Mlen$ is a weak semimeasure with $\Mlen(\emptyset)=\sum_iw_i\le1$ (each term is a product of probability vectors). On distinct inputs \cref{prop:setting:order} gives cumulative loss $-\log(\Mlen(D_k)/\Mlen(\emptyset))\le\log(1/w_i)+\sum_{j\le k}-\log q_i(y_j\mid x_j)$. If $|x_j|\ge n_A$, machine $i$ is active and $T_A(x_j)\le c_At(|x_j|)\le g(|x_j|)t(|x_j|)$, so $q_i(\cdot\mid x_j)=A(\cdot\mid x_j)$. If $|x_j|<n_A$, either $q_i(\cdot\mid x_j)=A(\cdot\mid x_j)$ (the machine is active and happens to finish in time) or $q_i(\cdot\mid x_j)=u$; in both cases $-\log q_i(y_j\mid x_j)\le\log|Y|-\log A(y_j\mid x_j)$, because $-\log A\ge0$. Subtract $A$'s loss. Distinctness bounds the number of inputs shorter than $n_A$ by $2^{n_A}-1$.
\end{proof}

\emph{Proof sketch of \cref{thm:comp:timebounded}(b).} Store $v_i:=w_i|Y|^k\prod_{j\le k}q_i(y_j\mid x_j)$ for the machines with $n_i\le\max_{j\le k}|x_j|$; every other machine has $v_i=w_i$. Then $\Mlen(D_k)=|Y|^{-k}\sum_iv_i$. A machine inactive at $x_k$ has $q_i=u$, and $|Y|u(y)=1$, so its $v_i$ does not change in round $k$; hence
\[ \Mlen(y\mid D_{k-1},x_k)=\frac{\sum_{i\text{ active at }x_k}v_iq_i(y\mid x_k)+\big(V-\sum_{i\text{ active at }x_k}v_i\big)/|Y|}{V},\qquad V:=\sum_iv_i, \]
and an update touches only the $\#\{i:n_i\le|x_k|\}\le\log\log|x_k|+1$ active weights, with $V$ maintained by adding their changes to its start value $\sum_iw_i$. \emph{Precision.} Let the learner store approximations $\hat v_i$ and use $\hat V_{k-1}$, computed exactly from them, with each active $\hat v_i$ replaced in round $k$ by $\hat v_i|Y|q_i(y_k\mid x_k)$ up to relative error $\varepsilon_k$. Its prediction $\hat p_k$ is a probability vector, and $\hat V_k\le(1+\varepsilon_k)|Y|\hat p_k(y_k)\hat V_{k-1}$, so the cumulative loss is at most $k\log|Y|+\log\hat V_0-\log\hat V_k+\sum_{j\le k}\log(1+\varepsilon_j)$. Since $\hat V_0\le1$ and $\hat V_k\ge\hat v_i\ge w_i|Y|^k\prod_{j\le k}q_i(y_j\mid x_j)(1-\varepsilon_j)$, the regret bound of (a) holds up to $\sum_j\log\frac{1+\varepsilon_j}{1-\varepsilon_j}=O(1)$ for $\varepsilon_j=j^{-2}$. The machine model (RAM versus multitape) and the cost of running the clocks are not written out; this part is a proof sketch.

\begin{remark}[size of the constant]\label{rem:comp:constant}
\status{proved}\src{context Rem 3.1a}
With $w_i=2^{-i}$ and $n_i=2^{2^i}$, the constant $\log(1/w_i)+(2^{n_A}-1)\log|Y|$ of \cref{thm:comp:timebounded}(a) is of order $2^{2^{2^i}}$, triple exponential in the index $i$, and $i$ can be exponential in the description length of $A$. Weights $w_A=2^{-\ell(A)}$ from a prefix-free code make $\log(1/w_A)$ the description length, but the activation term remains. H\"anni's sequence predictor starts machine $i$ at the round $\ell_1$ with $\log\log\ell_1\ge i$ \citep{hanni2026notes}, so his regret, $\max(\ell_1,\ell_2)$ bits, is doubly exponential in $i$. Length activation adds one exponential, because up to $2^{n_i}$ distinct inputs are shorter than $n_i$; it buys a per-round cost that depends only on the current input, as an x-only learner's should. The round-activated \cref{prop:comp:round} returns to the doubly exponential form, in $i$ and in $c_A$. A hybrid that activates machine $i$ at round $k_i$ but clocks by input length has, on distinct inputs, regret at most $\log(1/w_i)+(k_i-1)\log|Y|+\#\{j:|x_j|<n'_A\}\log|Y|$ with $n'_A:=\min\{n:g(n)\ge c_A\}$, by the proof of (a); it removes the extra exponential in $i$ but not in $c_A$, and its per-round cost grows slowly with the round count.
\end{remark}

\begin{proof}[Proof of \cref{prop:comp:history} (regret)]
$\tilde n_k$ and $t$ are nondecreasing, so for $k\ge k_0$ and $j\le k$ one round of $A$ at time $j$ costs at most $c_At(\tilde n_j)\le g(k)t(\tilde n_k)$. Hence in every learner round $k\ge k_0$ the budget $2g(k)t(\tilde n_k)$ completes at least two further rounds of $A$ (counting a partly done one), unless the replay reaches the present. If $r_k$ is the number of rounds of $A$ completed after learner round $k$, then $r_k\ge\min(k,r_{k-1}+2)$ for $k\ge k_0$, so $r_k\ge\min(k,2(k-k_0+1))=k$ for $k\ge2k_0-2$. A caught-up replay stays caught up, since one round costs at most the budget. So $q_i^{(k)}=A(\cdot\mid h_{k-1},x_k)$ for $k\ge2k_0-2$; for the earlier rounds $q_i^{(k)}\in\{A(\cdot\mid h_{k-1},x_k),u\}$ and $-\log q_i^{(k)}(y_k)\le\log|Y|-\log A(y_k\mid h_{k-1},x_k)$. Sum as in \cref{prop:comp:round}. If the budget depended only on $|x_k|$, a stream of short inputs after a long one would leave a budget too small to replay even that one round, and the catch-up could stall forever.
\end{proof}

\begin{proof}[Proof of \cref{prop:comp:occam}(b)]
For step $n+1$ with $n\ge1$ and every $\varepsilon$, $\Prob[\text{mistake}]=\Ex[\mathrm{err}_\mu(h_n)]\le\varepsilon+|H_f|(1-\varepsilon)^n$. With $\varepsilon_n:=\ln(|H_f|(n+1)^2)/n$ this is at most $\varepsilon_n+(n+1)^{-2}$. Summing over steps $1,\dots,n_0$ (step $1$ costs at most $1$), with $\ln(n+1)\le\ln n+1/n$ and $\sum_{n<n_0}\frac{\ln n}n\le\frac12\ln^2n_0+\frac1e$, gives at most $5.7+(\ln|H_f|)(1+\ln n_0)+\ln^2n_0$. After step $n_0$, each mistake is made by the current hypothesis, which lies in $H_f$, is consistent with the first $n_0$ samples, is wrong at a point of positive mass (so has error at least $p_{\min}$), and is eliminated by the mistake. The expected number of such hypotheses is at most $|H_f|(1-p_{\min})^{n_0}\le1$.
\end{proof}

\begin{proof}[Proof of \cref{thm:comp:circuitbayes}]
(a) $\sum_{s>S}2^{-s}=2^{-S}$. If no circuit of size at most $S$ agrees with $f$ on $B$, then $\Wcirc(f|_B)<2^{-S}$, strictly because at every size some constant circuit disagrees with $f$ on $B\ne\emptyset$; so $S<L_B$. Take $S=\lceil L_B\rceil$. Also $\Wcirc(f|_B)\ge w(C_{\min})=2^{-\csize_B(f)}/N_{\csize_B(f)}$, and $\log N_s\le s(4+2\log(d+1+s))$.

(b) First inequality: $P_i\le-\ln(1-P_i)$, and $\sum_i-\log(1-P_i)=L_B$ by \cref{prop:setting:order} with $\Wcirc(\emptyset)=1$. Second: fix $i$ and $S:=\lceil L_B\rceil$; let $\mathcal R$ and $\mathcal W$ be the circuits consistent with $D_{i-1}$ that give $x_i$ the right and the wrong label. The exception gadget at $x_i$ maps $\mathcal W_{\le S}$ (size at most $S$) injectively into $\mathcal R$, since earlier inputs differ from $x_i$, and $\Gamma$ is increasing; so $w(\mathcal W_{\le S})\le2^{\Gamma(S)}w(\mathcal R)$. Also $w(\mathcal W_{>S})\le2^{-S}\le2^{-L_B}=\Wcirc(f|_B)\le w(\mathcal R)$, because $\mathcal R$ contains every circuit that agrees with $f$ on $B$. So $P_i/(1-P_i)\le2^{\Gamma(S)}+1$ and $-\log(1-P_i)\le\Gamma(S)+2$. In $L_B=\sum_i-\log(1-P_i)$, terms with $P_i\le\frac12$ contribute at most $2P_i$ (the chord from $(0,0)$ to $(\frac12,1)$), and at most $2\mathcal E$ indices have $P_i>\frac12$. So $L_B\le2\mathcal E(\Gamma(S)+3)$.

Explicit form. $\gamma\le d+1$ and $d+1+S+\gamma\le2d+4+L_B$, so $L_B\le R(L_B)$ with $R(L):=2\mathcal E[(d+1)(5+2\log(2d+4+L))+3]$. $R$ is concave and increasing with $R(0)>0$, so $\{L\ge0:L\le R(L)\}$ is an interval $[0,L^*]$, and it suffices to show $K>R(K)$ for $K:=16a\Lambda$, where $a:=\mathcal E(d+1)$, $\Lambda:=X+Y$, $X:=\log(2(d+2))\ge\log6$, $Y:=\log(\mathcal E+1)$. Since $6\mathcal E\le3a$, $R(K)\le13a+4a\log(2d+4+K)$, and $2d+4+16a\Lambda\le2(d+2)(1+8\mathcal E\Lambda)\le2(d+2)(1+8\mathcal E)(1+X)(1+Y)$. So it suffices that
\[ \big[12X-4\log(1+X)-13\big]+\big[16Y-4\log(1+8\mathcal E)-4\log(1+Y)\big]>0. \]
The first bracket is increasing in $X$ and equals $10.65$ at $X=\log6$. For the second, $\log(1+Y)\le Y+0.087$ for all $Y\ge0$, and $\log(1+Y)\le Y$ for $Y\ge1$. If $\mathcal E\le1$, then $(1+8\mathcal E)/(1+\mathcal E)\le\frac92$ and the bracket is at least $12Y-4\log4.5-4(Y+0.087)\ge-9.03$. If $\mathcal E\ge1$, then $(1+8\mathcal E)/(1+\mathcal E)<8$ and the bracket exceeds $12Y-12-4Y\ge-4$. The sum is positive. A bisection on a grid ($\mathcal E\in[10^{-3},10^5]$, $d\in[1,10^4]$; code in the accompanying repository) gives the largest ratio $L^*/(\mathcal E(d+1)\log(2(\mathcal E+1)(d+2)))=10.85$, at $\mathcal E\approx0.63$, $d=1$.
\end{proof}

\begin{lemma}[M\"obius lemma]\label{lem:comp:mobius}
\status{proved; computed (context c6(a), $d\le4$)}\src{context Lemma 3.11a}
Let $1\le k\le d$ and $h:\{0,1\}^d\to\{0,1\}$ with $h(1^d)=1$ and $h(x)=0$ whenever $|x|<k$. Then (i) $h$ depends on at least $k$ variables; (ii) if it depends on exactly $k$ variables, then $h=\mathrm{AND}_S:=\bigwedge_{j\in S}x_j$ for some $|S|=k$.
\end{lemma}

\begin{proof}
Write $h=\bigoplus_Sc_S\prod_{j\in S}x_j$ over GF$(2)$. Evaluating at the indicator $1_T$ of $T$ gives $h(1_T)=\bigoplus_{S\subseteq T}c_S$, so $c_S=\bigoplus_{T\subseteq S}h(1_T)$. If $|S|<k$, every $T\subseteq S$ has $|T|<k$, so $c_S=0$. Since $h\ne0$, some $c_S=1$ with $|S|\ge k$; by uniqueness of this representation, $h$ depends on every variable of such a monomial, which gives (i). For (ii), let $h$ depend only on the variables in $R$, $|R|=k$. For $a\in\{0,1\}^R$ with $a\ne1_R$, the input equal to $a$ on $R$ and $0$ elsewhere has weight less than $k$, so $h=0$ there; and $h(1_R,\cdot)=h(1^d)=1$. So $h=\mathrm{AND}_R$.
\end{proof}

\begin{proof}[Proof of \cref{thm:comp:ladder}]
Stage $k$ ($1\le k\le d-1$) begins when the data are $h(1^d)=1$ and $h(x)=0$ for $|x|<k$, that is, after $1^d$, $0^d$ and all inputs of weight below $k$ have been shown with the labels of $f$.

\emph{Formulas.} A formula with $s$ leaves depends on at most $s$ variables. By \cref{lem:comp:mobius}(i) every function consistent at the start of stage $k$ has formula size at least $k$; one of size exactly $k$ depends on exactly $k$ variables and is a $k$-conjunction by \cref{lem:comp:mobius}(ii). Every $\mathrm{AND}_S$ with $|S|=k$ is consistent and has a formula of size $k$. So the consistent functions of least size are exactly the $\binom dk$ $k$-conjunctions, and Occam's hypothesis is one of them, since description length increases with size.

\emph{Circuits, $k\ge2$.} In a circuit with $s$ gates, let $t\le s$ be the number of gates from which the output can be reached. Each of the $t-1$ reachable non-output gates occupies at least one of the $2t$ input slots of reachable gates, so at most $t+1$ slots hold variables, and the function depends on at most $s+1$ variables. By \cref{lem:comp:mobius}(i) a consistent circuit has $s\ge k-1$. If $s=k-1$, then $t=s$ and exactly $k$ slots hold variables, which must be $k$ distinct ones; so the function depends on exactly $k$ variables and is a $k$-conjunction. $\mathrm{AND}_S$ has a circuit of $k-1$ gates. So again the least-size consistent functions are the $k$-conjunctions. (For $k=1$ this fails: the one-gate circuit $x_1\wedge x_2$ is consistent at the start of stage $1$.)

\emph{Every input of weight $k$ is a mistake.} For $|S|=|S'|=k$, $\mathrm{AND}_{S'}(1_S)=1$ iff $S'\subseteq S$ iff $S'=S$. If Occam's hypothesis in stage $k$ is $\mathrm{AND}_S$, then $1_S$ is unrevealed (else $\mathrm{AND}_S$ would be inconsistent), and the ordering shows it: Occam predicts $1$, the label is $f(1_S)=0$, and exactly $\mathrm{AND}_S$ among the $k$-conjunctions is eliminated. So all $\binom dk$ inputs of weight $k$ are mistakes, after which stage $k+1$ begins.

\emph{Count and description length.} Formulas: $\sum_{k=1}^{d-1}\binom dk=2^d-2$; circuits: $\sum_{k=2}^{d-1}\binom dk=2^d-d-2$. At $k=d$ the argument shows that $f$ has least size $d$ (formulas) or $d-1$ (circuits), so $\ell(f)=d+\log(\mathrm{Cat}(d-1)\,3^{d-1}(2d+2)^d)$ for formulas, with $\mathrm{Cat}$ the Catalan numbers counting tree shapes, and $\ell(f)=(d-1)+\log N_{d-1}$ for circuits. Both are $\Theta(d\log d)$: $\log((2d+2)^d)$ dominates $\mathrm{Cat}(d-1)\le4^d$, and $\log N_{d-1}=\sum_{j<d}(4+2\log(d+1+j))$. So $d=\Theta(\ell/\log\ell)$ and the number of mistakes is at least $2^{d-1}=2^{\Omega(\ell/\log\ell)}$ for $d\ge4$. The trivial upper bound is $\min(2^d,|H_f|-1)\le2^{\ell+1}$.
\end{proof}

\begin{remark}[a second adversary]\label{rem:comp:adversary}
\status{proved; computed (context c11)}\src{context Rem 3.11b}
For circuits the count $2^d-d-2$ also follows without the stage lemma, for every tie rule. Reveal $1^d$; then, while Occam's hypothesis $h$ has $\csize(h)<d-1$, let $V_h$ be the set of variables $h$ depends on and reveal $1_{V_h}$. Each step is a mistake: by the slot count, $|V_h|\le\csize(h)+1<d$; $h$ depends only on $V_h$, so $h(1_{V_h})=h(1^d)=1$, while $f(1_{V_h})=0$; and the point is new, or $h$ would be inconsistent. The phase reveals a new point at each step, so it ends, and then $\csize(h)\ge d-1$. Fix $2\le k\le d-1$ and $|T|=k$, so $\csize(\mathrm{AND}_T)=k-1$. While $\mathrm{AND}_T$ is consistent, Occam's hypothesis precedes it, so $\csize(h)\le k-1<d-1$: the phase continues and reveals points of weight $|V_h|\le k$. When the phase ends $\mathrm{AND}_T$ has been killed by a revealed point $x$ with $\mathrm{AND}_T(x)=1\ne f(x)$, revealed while $\mathrm{AND}_T$ was consistent; so the ones of $x$ contain $T$, $x\ne1^d$ and $|x|\le k$, hence $x=1_T$, revealed as a mistake. So every point of weight $2,\dots,d-1$ is a mistake. A direct computation under seven tie rules ($d=3,4,5$) confirms this. (For formulas the slot count gives $|V_h|\le\csize(h)$, and every point of weight $1,\dots,d-1$ is a mistake.)
\end{remark}

A direct computation (code in the accompanying repository) checks \cref{lem:comp:mobius} for all functions with $d\le4$, the least-size statement for formulas and circuits with $d=3,4,5$, and simulates Occam on the ladder: $7,15,31$ mistakes (formulas) and $4,11,26$ (circuits) for $d=3,4,5$, with ties broken by truth-table index. Under seven tie rules every input of every level the theorem covers is a mistake. The ratio of $\log(\text{mistakes})$ to $\ell/\log\ell$ stays between $0.64$ and $0.82$ (formulas) and between $0.35$ and $0.59$ (circuits) for $d=3,\dots,40$.
